\subsection{Event Definition}\label{sec:EVENT_DEFINITION}
\ac{DUNE}'s solar analysis workflow starts with calibrating the clustered hits from the \ac{CC} interaction of simulated neutrino events with the MARLEY generator. The \ac{CC} channel is chosen to be the default based on its dominant cross--section in the low--energy range, as well as its richer topology. All energy and spacial resolution results presented here should be considered and underestimation (or at least compatible) to the ones expected for the \ac{ES} channel. Each cluster's information is reduced to a point--like object characterized by its number of hits (\texttt{nhits}), charge (summed over the charge of the individual hits) and a 2.5D--reconstructed vertex $<t, y, z>$ computed from the charge--weighted average of the individual hits.

\begin{table}[ht]
    \centering
    \begin{minipage}{0.45\linewidth}
        \centering
        \input{Tables/hd_1x2x6_centralAPA_marley_official_Signal_AdjClusters.txt}
        \caption{\ac{HD} average number of adjacent clusters per primary signal cluster as a function of their 3D--distance for the lateralAPA background scenario.\label{tab:ADJCLUSTERNUM_RADIAL_HD}} 
    \end{minipage}\hfill
    \begin{minipage}{0.45\linewidth}
        \centering
        \input{Tables/vd_1x8x14_3view_30deg_nominal_marley_official_Signal_AdjClusters.txt}
        \caption{\ac{VD} average number of adjacent clusters per primary cluster as a function of their 3D--distance for the nominal background scenario.\label{tab:ADJCLUSTERNUM_RADIAL_VD}} 
    \end{minipage}
\end{table}

Clusters are associated to others based on their relative charges and radial distance. For the radial distance comparison between clusters, the difference in their arrival times is converted to the missing spacial coordinate by taking into account the known drift velocity in \ac{LAr}. Notice that this computation only holds true for correlated energy deposits (e.g. nuclear de--exitation gammas from a \ac{CC} neutrino interaction in \ac{LAr} following the emission of a main electron). In the case of non-correlated background events occurring throughout the full detector volume, their distance along the drift coordinate with respect to other clusters cannot be truthfully determined. From a reconstruction point of view, this leads to the presence of a constant background floor in the vicinity of any signal cluster that corresponds to the projected activity in the interacting drift volume (see \cref{tab:ADJCLUSTERNUM_RADIAL_HD} and \cref{tab:ADJCLUSTERNUM_RADIAL_VD}). These clusters are divided into:

\begin{itemize}
    \item \textbf{primary} with charge bigger than any other's cluster charge up to \SI{1}{\meter} radial distance.
    \item \textbf{adjacent} with charge smaller than the one of a primary cluster closer than \SI{1}{\meter} radial distance.
\end{itemize}

We define a low--energy event as the combination of a primary cluster together with its adjacent clusters and a list of flashes that are matching candidates for the drift coordinate reconstruction (see \ref{sec:FLASH_MATCHING}). In any neutrino analysis it will become important to try to distinguish (at least on average) the adjacent cluster distributions for any low--energy event, not only to distinguish background activity (e.g. neutron captures) from signal events, but also to identify the different interaction channels of the neutrino interaction (\ac{CC} events produce de--exitation gammas that are not present in \ac{ES} interactions). Consequently, the higher the projected background activity in any of the isolated drift volumes of \ac{DUNE}'s \ac{FD}s, the harder particle identification and channel separation will become. Notice that there is nothing preventing a single neutrino interaction from producing more than one low--energy event. Rarely secondary gammas can scatter further than \SI{1}{\meter} in \ac{LAr} or a background--induced miss--match in the \ac{TPC} reconstruction could lead to blips appearing in uncorrelated regions of the detector.

\begin{figure}[ht!]
    \centering
    \includegraphics[width=0.49\linewidth]{Images/hd_1x2x6_marley_official_PreselectionE_Efficiency.png}
    \includegraphics[width=0.49\linewidth]{Images/vd_1x8x14_3view_30deg_marley_official_PreselectionE_Efficiency.png}
    \caption{\ac{HD} (left) and \ac{VD} (right) detection efficiency for signal clustering parameters of 3 adjacent channels and \SI{12.5}{\micro\second}.\label{fig:DETECTION_EFFICIENCY}}
\end{figure}

For the purpose of this work, secondary low--energy events coming from a single neutrino interaction will always be ignored. This is to ensure that the reconstruction performance is evaluated for the most complete reconstruction object in any interaction but in any physics analysis these should be kept and treated as part of the signal or even background.
